\chapter{Sprint 2: Psychologist Onboarding System}
\label{chap:sprint2-onboarding}

\section*{Introduction}
\addcontentsline{toc}{section}{Introduction}

Sprint~2 focuses on the design and implementation of a structured onboarding
pipeline for psychologists. In a mental-health platform, onboarding is more
than a registration form: it is a security boundary and a trust mechanism. The
platform must ensure that only verified and qualified professionals are able to
offer services, appear in search results, and access sensitive patient-facing
features.

This sprint introduces an end-to-end onboarding workflow covering: (i) profile
creation and completion, (ii) secure submission of professional credentials,
(iii) AI-assisted pre-screening to detect obvious inconsistencies or missing
elements, and (iv) an administrative review process that produces a final
approve/reject decision. A key requirement is auditability: every important
transition (submission, verification, admin decision) is recorded with
timestamps and actor identity to guarantee traceability and accountability.

\section{Sprint Objectives and Scope}
\label{sec:sprint2-objectives}

The main objective of Sprint~2 is to deliver a complete onboarding system that
transitions a psychologist from an initial account state to an \emph{Active}
state only after verification. The scope includes both user-facing components
and internal operational tooling:

\begin{itemize}
    \item Collect comprehensive professional information (identity, practice
    details, specialisations, and experience).
    \item Enable secure upload and management of credential documents.
    \item Integrate an AI-assisted verification step that produces a structured
    result (risk score, extracted entities, and flags).
    \item Provide an admin review queue and a decision workflow with explicit
    reasoning and audit logging.
    \item Enforce access control so that only approved psychologists can use
    privileged features and become discoverable on the platform.
\end{itemize}

\textbf{Out of scope.} Sprint~2 does not attempt to fully automate professional
licence validation against external government registries (where such APIs are
often unavailable or inconsistent). Instead, the system is designed so that AI
pre-screening supports (but does not replace) human administrative judgement.

\textbf{Success Criteria.}
\begin{itemize}
    \item A psychologist can complete onboarding without admin intervention
    until the final approval step.
    \item Admin decisions are fully traceable (who, when, what, and why).
    \item The onboarding status model has clear state transitions and prevents
    invalid moves (e.g., approving an incomplete submission).
    \item Uploaded credentials are stored securely and only accessible to
    authorised roles.
\end{itemize}

\section{Sprint Planning}
\label{sec:sprint2-planning}

\subsection{Assumptions}
\begin{itemize}
    \item Sprint duration: 2 weeks.
    \item Agile team structure: Product Owner, backend developer(s), frontend
    developer(s), QA, and an admin representative for workflow validation.
    \item Security and privacy constraints are treated as first-class
    requirements (least privilege, encrypted storage, and audit logging).
\end{itemize}

\subsection{Deliverables}
\begin{itemize}
    \item Onboarding UI (profile form, document checklist, submission and status
    tracking).
    \item Backend APIs for onboarding records, document upload metadata, and
    status transitions.
    \item AI verification pipeline integration (job submission, result
    persistence, and failure handling).
    \item Admin dashboard (review queue, submission detail view, decision
    actions, and audit log view).
    \item Testing pipeline covering validation, permissions, and end-to-end
    onboarding scenarios.
\end{itemize}

\section{Stakeholders and Roles}
\label{sec:sprint2-stakeholders}

Sprint~2 formalises responsibilities across three main actors:

\begin{itemize}
    \item \textbf{Psychologist:} provides profile data and credentials, submits
    onboarding request, and monitors status updates and feedback.
    \item \textbf{Admin:} reviews submissions, interprets AI verification
    results, and approves or rejects onboarding with recorded reasoning.
    \item \textbf{AI Verification Service:} analyses submitted documents and
    returns structured evidence to support admin decisions (risk score, flags,
    and extracted fields).
\end{itemize}

\section{Sprint Backlog}
\label{sec:sprint2-backlog}

\renewcommand{\arraystretch}{1.4}

\subsection{User Stories (Onboarding Flow)}
\begin{table}[H]
\centering
\begin{tabularx}{\textwidth}{|c|X|X|}
\hline
\rowcolor{medgray}
\textbf{ID} & \textbf{Description} & \textbf{Acceptance Criteria} \\
\hline

US-2.1 & Creation of professional profile &
\begin{itemize}[leftmargin=*]
\item Validate required fields (identity, contact, location, specialisation).
\item Allow saving as draft and resuming later.
\item Persist a clear ``completion'' indicator to guide the user.
\end{itemize} \\ \hline

US-2.2 & Upload credential documents &
\begin{itemize}[leftmargin=*]
\item Validate file type and size; reject unsupported uploads.
\item Categorise documents (licence, degree, ID, other).
\item Confirm secure storage and display upload status per document.
\end{itemize} \\ \hline

US-2.3 & Submit onboarding request &
\begin{itemize}[leftmargin=*]
\item Block submission if required fields/documents are missing.
\item Assign onboarding status and store submission timestamp.
\item Trigger AI verification and store a job reference.
\end{itemize} \\ \hline

US-2.4 & Track onboarding status &
\begin{itemize}[leftmargin=*]
\item Display status and next steps clearly (what is pending and why).
\item If rejected, show a human-readable rejection reason.
\item Provide guidance for resubmission when applicable.
\end{itemize} \\ \hline

\end{tabularx}
\caption{User Stories: Onboarding Flow}
\end{table}

\subsection{AI and Admin Workflow}
\begin{table}[H]
\centering
\begin{tabularx}{\textwidth}{|c|X|X|}
\hline
\rowcolor{medgray}
\textbf{ID} & \textbf{Feature} & \textbf{Description} \\
\hline

US-2.5 & AI Verification & AI analyses documents and returns risk score, extracted fields, and flags. \\ \hline
US-2.6 & Admin Queue & Displays pending onboarding requests with filtering and prioritisation. \\ \hline
US-2.7 & Admin Decision & Admin approves/rejects with mandatory logging, reason, and timestamp. \\ \hline

\end{tabularx}
\caption{User Stories: AI and Admin Workflow}
\end{table}

\subsection{Technical Tasks}
\begin{table}[H]
\centering
\begin{tabularx}{\textwidth}{|X|X|}
\hline
\rowcolor{medgray}
\textbf{Task} & \textbf{Description} \\
\hline

State Machine & Define onboarding states, transitions, and invariants. \\ \hline
Secure Storage & Implement encrypted credential storage and controlled access paths. \\ \hline
AI Integration & Integrate verification pipeline (async jobs, retries, persistence). \\ \hline
Admin Dashboard & Build moderation interface and decision endpoints with audit trail. \\ \hline
Testing & Add unit, integration, and end-to-end tests for onboarding scenarios. \\ \hline

\end{tabularx}
\caption{Technical Tasks}
\end{table}

\section{Functional Specification}
\label{sec:sprint2-functional}

\subsection{Profile Data Collection}
The onboarding profile is designed as a structured record with both required
and optional fields. Required fields focus on identity and professional
credibility; optional fields enhance matching and discovery:

\begin{itemize}
    \item \textbf{Identity and contact:} full name, verified email/phone,
    country/city, profile photo (optional, subject to moderation policy).
    \item \textbf{Professional details:} licence/registration number (when
    applicable), issuing authority, specialisations, spoken languages, years of
    experience, and short professional biography.
    \item \textbf{Operational preferences:} availability windows, online/offline
    capability, consultation types offered, and pricing policy (if the platform
    supports it).
\end{itemize}

Client-side and server-side validation rules are aligned to prevent inconsistent
states and to ensure the backend remains the source of truth. Draft saving is
supported to reduce onboarding abandonment and allow incremental completion.

\subsection{Credential Submission and Document Handling}
Credential submission is implemented as a checklist-driven flow. Each document
upload creates a record containing:

\begin{itemize}
    \item Document type (licence, degree, ID, other).
    \item Original filename, MIME type, file size, and upload timestamp.
    \item Storage reference (e.g., encrypted object key) and integrity metadata
    (e.g., checksum) when supported.
    \item Visibility constraints enforcing that only the psychologist (for their
    own submission) and admins can access documents.
\end{itemize}

To maintain privacy and limit exposure of sensitive documents, the platform
should prefer short-lived, role-scoped access mechanisms (e.g., signed URLs or
streaming through the backend) rather than permanent public links.

\subsection{Submission, Locking, and Resubmission Policy}
Upon submission, the system validates completeness and transitions the record
into an immutable review state (policy-dependent). A common approach is to lock
critical fields and documents to avoid post-submission tampering; any changes
require explicit resubmission and generate a new audit event.

If rejected, the record includes an admin-provided reason and can either:
(i) allow edits and resubmission while preserving the historical audit trail, or
(ii) require a new onboarding attempt linked to the same user. Sprint~2 targets
the first option because it minimises friction while preserving traceability.

\section{AI-Assisted Verification}
\label{sec:sprint2-ai}

The AI verification step is designed as a pre-screening mechanism that provides
evidence to admins. Typical checks include:

\begin{itemize}
    \item \textbf{OCR and entity extraction:} extract name, licence number,
    issuing body, and validity dates.
    \item \textbf{Consistency checks:} compare extracted entities against
    profile fields (e.g., name match, country/authority plausibility).
    \item \textbf{Completeness checks:} detect missing required document types
    or unreadable uploads (low OCR confidence).
    \item \textbf{Risk scoring:} compute an overall risk score with supporting
    flags to guide review prioritisation.
\end{itemize}

Because AI systems can produce false positives and false negatives, the output
is treated as \emph{advisory} rather than determinative. The admin interface
must clearly separate factual extracted fields from risk heuristics and provide
enough context to justify the final decision.

\section{Administrative Review and Auditability}
\label{sec:sprint2-admin}

\subsection{Admin Review Queue}
The admin queue provides operational visibility and prioritisation. Essential
capabilities include filtering by status (e.g., submitted, pending review,
flagged), sorting by submission date, and quickly opening a detailed view.

\subsection{Decision Workflow}
Admin decisions must be explicit and recorded. Approvals transition the
psychologist account to an \emph{Active} state and enable platform privileges.
Rejections require a reason and transition the record to \emph{Rejected}. The
system records:

\begin{itemize}
    \item Admin identity and timestamp.
    \item Decision type (approve/reject).
    \item Rationale text and, when applicable, references to AI flags or
    inconsistencies.
\end{itemize}

\subsection{Audit Log}
An audit log is maintained for all state transitions and security-sensitive
events (document upload, submission, AI result received, admin decision,
resubmission). This log supports operational debugging, dispute resolution, and
compliance requirements.

\section{Onboarding State Machine}
\label{sec:sprint2-state-machine}

Sprint~2 enforces onboarding progression through a controlled state machine.
The exact naming may vary in implementation, but the conceptual states are
stable:

\begin{table}[H]
\centering
\begin{tabularx}{\textwidth}{|l|X|X|}
\hline
\rowcolor{medgray}
\textbf{State} & \textbf{Meaning} & \textbf{Typical Transitions} \\
\hline

Draft & Profile and/or documents are being prepared. &
Draft $\rightarrow$ Submitted (after completeness validation). \\ \hline

Submitted & Submission is complete and awaiting verification. &
Submitted $\rightarrow$ Verifying; Submitted $\rightarrow$ Rejected (admin). \\ \hline

Verifying & AI verification job is running or pending results. &
Verifying $\rightarrow$ Pending Admin Review; Verifying $\rightarrow$ Verification Failed (retry). \\ \hline

Pending Admin Review & Admin review is required before activation. &
Pending Admin Review $\rightarrow$ Approved; Pending Admin Review $\rightarrow$ Rejected. \\ \hline

Approved & Onboarding accepted; psychologist can be activated. &
Approved $\rightarrow$ Active (system action). \\ \hline

Rejected & Onboarding rejected; edits and resubmission allowed. &
Rejected $\rightarrow$ Draft (edit); Rejected $\rightarrow$ Submitted (resubmit). \\ \hline

\end{tabularx}
\caption{Onboarding state machine (conceptual model)}
\end{table}

The state machine prevents invalid operations such as approving without a
submission, submitting without required documents, or exposing privileged
features before activation.

\section{System Design}
\label{sec:sprint2-design}

\subsection{Use Case Diagram}
\begin{figure}[H]
\centering
\fbox{\parbox{0.85\textwidth}{\centering Placeholder for Use Case Diagram (Sprint~2)}}
\caption{Use Case Diagram: Psychologist Onboarding}
\end{figure}

\subsection{Class Diagram}
\begin{figure}[H]
\centering
\fbox{\parbox{0.85\textwidth}{\centering Placeholder for Class Diagram (Onboarding, Documents, AuditLog)}}
\caption{Class Diagram: Onboarding Domain}
\end{figure}

\subsection{Sequence Diagram}
\begin{figure}[H]
\centering
\fbox{\parbox{0.85\textwidth}{\centering Placeholder for Sequence Diagram (Submit $\rightarrow$ AI Verify $\rightarrow$ Admin Decision)}}
\caption{Sequence Diagram: End-to-End Onboarding}
\end{figure}

\subsection{State Machine Diagram}
\begin{figure}[H]
\centering
\fbox{\parbox{0.85\textwidth}{\centering Placeholder for State Machine Diagram}}
\caption{Onboarding State Machine}
\end{figure}

\section{Testing and Validation Strategy}
\label{sec:sprint2-testing}

\subsection{Unit Testing}
\begin{itemize}
    \item Validate profile field rules and error messages.
    \item Test state transition guards (valid vs.\ invalid moves).
    \item Verify RBAC permissions for psychologist vs.\ admin operations.
\end{itemize}

\subsection{Integration Testing}
\begin{itemize}
    \item API communication between client and backend for onboarding actions.
    \item Document storage interactions (upload metadata, access controls).
    \item AI pipeline integration (job creation, result persistence, retries).
\end{itemize}

\subsection{End-to-End Testing}
\begin{itemize}
    \item Complete onboarding workflow from draft to admin approval.
    \item Admin review flow with both approve and reject outcomes.
    \item Access control checks ensuring only approved psychologists become
    discoverable and gain privileged access.
\end{itemize}

\subsection{Test Environment}
Testing is conducted in a staging environment configured with production-like
security controls and monitoring. Where external AI services are involved,
tests use deterministic fixtures or mocked responses to ensure reliability.

\section{Sprint Review and Retrospective}
\label{sec:sprint2-review}

\subsection{Sprint Review}
\begin{itemize}
    \item Demonstration of the onboarding UI (drafting, uploading, submission).
    \item Demonstration of AI verification results integrated into admin review.
    \item Admin decision workflow including audit trail visibility.
\end{itemize}

\subsection{Retrospective}
\textbf{Challenges:}
\begin{itemize}
    \item AI integration complexity (latency, failures, and confidence handling).
    \item Security constraints for document access and storage.
    \item Balancing user experience with strict validation requirements.
\end{itemize}

\textbf{Improvements:}
\begin{itemize}
    \item Improve user feedback (clearer checklist and missing-item guidance).
    \item Optimise AI processing (async jobs, retries, and better status UX).
    \item Extend admin tooling (bulk actions, prioritisation, and analytics).
\end{itemize}

\section{Conclusion}
\label{sec:sprint2-conclusion}

Sprint~2 delivers a robust onboarding pipeline that combines secure credential
handling, AI-assisted pre-screening, and an accountable administrative approval
workflow. By enforcing a controlled state machine and a complete audit trail,
the platform establishes the trust and governance needed to safely connect
patients with verified psychologists.

